% MCM/ICM (美赛) LaTeX Template
% Math Modeling Contest Agent
\documentclass[12pt]{article}

% === Packages ===
\usepackage[utf8]{inputenc}
\usepackage{amsmath,amssymb,amsfonts}
\usepackage{graphicx}
\usepackage{booktabs}
\usepackage{multirow}
\usepackage{longtable}
\usepackage{enumitem}
\usepackage{hyperref}
\usepackage[margin=1in]{geometry}
\usepackage{fancyhdr}
\usepackage{float}
\usepackage{caption}
\usepackage{subcaption}
\usepackage{algorithm}
\usepackage{algpseudocode}
\usepackage{xcolor}
\usepackage{listings}
\usepackage{setspace}
\usepackage{titling}

% === Page Setup ===
\geometry{left=2.5cm,right=2.5cm,top=2.5cm,bottom=2.5cm}
\onehalfspacing

% === Header/Footer ===
\pagestyle{fancy}
\fancyhf{}
\fancyhead[L]{\small Team \#XXXXXX}
\fancyhead[R]{\small Page \thepage}
\renewcommand{\headrulewidth}{0.4pt}

% === Code Style ===
\lstset{
    basicstyle=\ttfamily\small,
    breaklines=true,
    frame=single,
    numbers=left,
    numberstyle=\tiny,
    xleftmargin=2em,
}

% === Title ===
\pretitle{\begin{center}\LARGE\bfseries}
\posttitle{\end{center}}
\title{Problem Chosen: [A/B/C/D/E/F]\\[0.5em]\large Paper Title}
\author{}
\date{}

\begin{document}

\maketitle
\thispagestyle{empty}

% === Summary ===
\begin{center}
\textbf{\Large Summary}
\end{center}

% TODO: Summary written by Agent
This paper addresses [problem background] by developing [model type] models.

\textbf{Keywords:} keyword1; keyword2; keyword3; keyword4

\newpage
\tableofcontents
\newpage

% ==========================================
% 1. Introduction
% ==========================================
\section{Introduction}

\subsection{Problem Background}
% TODO: Agent fills in

\subsection{Restatement of the Problem}
% TODO: Agent fills in

\subsection{Our Work}
% Overview of approach

% ==========================================
% 2. Assumptions and Justifications
% ==========================================
\section{Assumptions and Justifications}

\begin{itemize}
    \item \textbf{Assumption 1:} Description. \textit{Justification:} Reason.
    \item \textbf{Assumption 2:} Description. \textit{Justification:} Reason.
\end{itemize}

% ==========================================
% 3. Notations
% ==========================================
\section{Notations}

\begin{table}[H]
\centering
\caption{Notation Table}
\begin{tabular}{cll}
\toprule
\textbf{Symbol} & \textbf{Description} & \textbf{Unit} \\
\midrule
$x$ & Variable description & - \\
\bottomrule
\end{tabular}
\end{table}

% ==========================================
% 4. Models
% ==========================================
\section{Models}

\subsection{Task 1: Model Name}
\subsubsection{Model Formulation}
% TODO: Agent fills in mathematical model

\subsubsection{Solution Approach}
% TODO: Agent fills in solution method

\subsubsection{Results and Analysis}
% TODO: Agent fills in results

\subsection{Task 2: Model Name}
% Same structure as above

\subsection{Task 3: Model Name}
% Same structure as above

% ==========================================
% 5. Sensitivity Analysis
% ==========================================
\section{Sensitivity Analysis}
% TODO: Agent fills in

% ==========================================
% 6. Model Evaluation
% ==========================================
\section{Model Evaluation}

\subsection{Strengths}
\begin{enumerate}
    \item Strength 1
    \item Strength 2
\end{enumerate}

\subsection{Weaknesses}
\begin{enumerate}
    \item Weakness 1
    \item Weakness 2
\end{enumerate}

\subsection{Future Work}
% TODO: Agent fills in

% ==========================================
% 7. Conclusion
% ==========================================
\section{Conclusion}
% TODO: Agent fills in

% ==========================================
% References
% ==========================================
\begin{thebibliography}{99}

\bibitem{ref1} Author. Title[J]. Journal, Year, Volume(Issue): Pages.
\bibitem{ref2} Author. Title[M]. Publisher, Year.

\end{thebibliography}

% ==========================================
% Appendices
% ==========================================
\appendix
\section{Code Listing}
% Key solver code

\end{document}
